\documentclass[12pt,a4paper]{report}

% --- PACKAGES ---
\usepackage[french]{babel}
\usepackage[T1]{fontenc}
\usepackage[utf8]{inputenc}
\usepackage{listings}
\usepackage{xcolor}
\usepackage[most]{tcolorbox}
\usepackage{geometry}

\geometry{margin=2.5cm}

% --- COULEURS ---
\definecolor{dialog-purple}{rgb}{0.5, 0.0, 0.5}
\definecolor{code-bg}{rgb}{0.98, 0.98, 0.98}

% --- STYLE DES BOITES (Design corrigé pour Yahya) ---
\newtcolorbox{dialogbox}[1]{
    enhanced,
    attach boxed title to top left={yshift=-2mm, xshift=2mm},
    colback=white,
    colframe=dialog-purple,
    fonttitle=\bfseries,
    colbacktitle=dialog-purple,
    title=#1,
    boxed title style={size=small, sharp corners},
    drop shadow,
    bottom=2mm
}

% --- CONFIGURATION C ---
\lstset{
    language=C,
    backgroundcolor=\color{code-bg},
    basicstyle=\ttfamily\footnotesize,
    keywordstyle=\color{blue},
    commentstyle=\color{gray},
    stringstyle=\color{purple},
    breaklines=true,
    frame=none,
    numbers=left,
    numberstyle=\tiny\color{gray}
}

\begin{document}

\chapter{Module de Boîtes de Dialogue (Dialogs)}

\section{Introduction}
Les boîtes de dialogue sont essentielles pour confirmer des actions critiques (comme la suppression d'un trajet dans le \textit{Smart Transport Scheduler}). Ce module utilise l'API moderne de GTK4 pour afficher des alertes sans geler l'interface utilisateur.

\section{Structures de Configuration}

\begin{dialogbox}{Structure : alert\_dialog\_config}
Cette structure contient tout le nécessaire pour construire une alerte : le message principal, les détails secondaires, et surtout un tableau de boutons personnalisables.
\end{dialogbox}

\begin{lstlisting}
typedef struct {
    GtkWindow *parent;
    const char *message;
    const char *detail;
    const alert_button_config *buttons;
    size_t button_count;
    int default_button;
    int cancel_button;
    bool modal;
    void (*on_response)(int response_index, gpointer user_data);
    gpointer user_data;
} alert_dialog_config;
\end{lstlisting}

\section{Macros et Valeurs par Défaut}

\begin{dialogbox}{Macro : ALERT\_DIALOG\_CONFIG\_DEFAULT}
Pour simplifier la vie du développeur, cette macro initialise une alerte "vide" mais sécurisée. Par défaut, la fenêtre est modale (elle bloque l'interaction avec la fenêtre parente jusqu'à sa fermeture).
\end{dialogbox}

\begin{lstlisting}
#define ALERT_DIALOG_CONFIG_DEFAULT ((alert_dialog_config){ \
    .parent = NULL, \
    .message = NULL, \
    .modal = true, \
    .default_button = -1, \
    .cancel_button = -1, \
    // ... initialisation du reste a zero ...
})
\end{lstlisting}

\section{Gestion de l'Asynchronisme}

\begin{dialogbox}{Le Contexte : alert\_async\_context}
Comme GTK4 fonctionne de manière asynchrone, nous devons créer une petite structure "tampon" pour garder en mémoire les informations de l'utilisateur pendant qu'il choisit son bouton.
\end{dialogbox}

\begin{lstlisting}
typedef struct {
    GtkAlertDialog *dialog;
    void (*on_response)(int response_index, gpointer user_data);
    gpointer user_data;
} alert_async_context;
\end{lstlisting}

\section{Implémentation de l'Affichage}

\begin{dialogbox}{Fonction : show\_alert\_dialog}
C'est la fonction principale. Elle transforme notre configuration en un objet \texttt{GtkAlertDialog}. Si une fonction \texttt{on\_response} est fournie, elle prépare le rappel (callback) pour traiter le clic de l'utilisateur.
\end{dialogbox}

\begin{lstlisting}
void show_alert_dialog(const alert_dialog_config *config) {
    GtkAlertDialog *dialog;
    dialog = gtk_alert_dialog_new("%s", config->message);
    
    // Configuration des boutons et du focus
    if (default_button >= 0) {
        gtk_alert_dialog_set_default_button(dialog, default_button);
    }

    // Lancement de l'alerte
    if (config->on_response != NULL) {
        // Mode asynchrone avec callback
        gtk_alert_dialog_choose(dialog, config->parent, NULL, on_alert_response, ctx);
    } else {
        // Simple affichage informatif
        gtk_alert_dialog_show(dialog, config->parent);
    }
}
\end{lstlisting}

\end{document}